% ===================================================================
%  કોષ્ટક નં. ૧૦ — દરિયો. બંદર.
%
%  Traduction, faite en 2026, en gujarati d'aujourd'hui, sur l'ido de
%  texto/io. Regles de la colonne : voir 10-kostak-01.tex. Une seule
%  scene, une seule numerotation. 37 blocs, 98 renvois.
%
%  C'EST LE TABLEAU OU CETTE COLONNE EST CHEZ ELLE, ET IL FALLAIT BIEN
%  QU'IL ARRIVE. Le Gujarat est une cote — Surat, Bharuch, Mandvi,
%  Porbandar, Veraval —, et le releve vient de lui faire traverser
%  neuf tableaux ou il empruntait a chaque page : le sapin, le seigle,
%  la chataigne, le coquelicot, la pervenche, le hetre, la faux. Ici,
%  rien. વહાણ le navire, સઢ la voile, ડોલકાઠી le mat, હલેસું la rame,
%  સુકાન le gouvernail, તૂતક le pont, લંગર l'ancre, ધક્કો le wharf,
%  ગોદી le dock, વખાર l'entrepot, જકાત la douane, ખલાસી le matelot,
%  દીવાદાંડી le phare, ભરતી et ઓટ le flux et le reflux, ભૂશિર le cap,
%  દ્વીપકલ્પ la presqu'ile, સામુદ્રધુની le detroit : tout est la, tout
%  est vieux, rien n'est decrit. Aucune autre colonne du releve n'a eu
%  un tableau entier sans un seul emprunt de matiere.
%
%  ET LE PORT DE GUERRE, LUI, EST ENTIEREMENT ANGLAIS. ક્રૂઝર,
%  ફ્રિગેટ, સબમરીન, ટૉર્પિડો, બોઇલર, ક્રેન, લેફ્ટનન્ટ, એડમિરલ,
%  ક્વાર્ટરમાસ્ટર. La coupure passe exactement entre les deux ports
%  du tableau, et elle a une date : cette cote a commerce par mer
%  pendant deux mille ans et n'a plus eu de marine de guerre a elle
%  apres l'arrivee des Portugais. UNE LANGUE GARDE LES MOTS DE CE
%  QU'ELLE A FAIT ET EMPRUNTE CEUX DE CE QU'ON LUI A FAIT. Le tableau
%  9 le disait des arbres — on garde le nom de ce qu'on a manie ; ce
%  tableau-ci le dit des metiers, ce qui est plus dur.
%
%  « ઘાટ » REVIENT, ET CE N'EST PAS UNE COINCIDENCE. Au tableau 9 il
%  rendait le defile, la coupure par ou l'on franchit une chaine ; ici
%  il rend le quai (58, 45). Le mot ne veut dire ni montagne ni port :
%  il veut dire l'endroit amenage par ou l'on passe d'un element dans
%  un autre — l'escalier qui descend au fleuve, la breche qui descend
%  de la crete, la pierre qui descend a la mer. Deux tableaux, deux
%  renvois, un seul mot, et il est juste les deux fois.
%
%  LE CHANTIER NAVAL (18) EST « વહાણવાડો », ET IL EXISTE ENCORE. On
%  construit toujours des boutres de bois a Mandvi, au Kutch, a la
%  main et sans plan, par des charpentiers qui tiennent le metier de
%  leur pere. Le mot n'est donc pas un mot d'archive : c'est le nom
%  d'un endroit ou l'on peut aller aujourd'hui. Et « સુકાની », celui
%  qui tient le સુકાન, est le mot ordinaire pour un chef — la langue
%  a tire son mot de dirigeant du gouvernail d'un navire, comme le
%  francais l'a tire du gouvernail d'un char.
%
%  UN SEUL TROU, ET IL EST INSTRUCTIF : LA QUILLE (29). Une cote qui a
%  vingt mots pour un navire n'en a pas un seul, en gujarati standard,
%  pour la piece de bois qui court sous la coque. Les charpentiers de
%  Mandvi en ont un, mais c'est un mot de chantier, pas un mot de
%  langue, et le releve ne s'en sert pas : on ecrit « તળિયાનો પાટડો »,
%  la poutre du fond. UNE LANGUE PEUT MANQUER PRECISEMENT LA OU ELLE
%  SAIT LE PLUS — parce que ce qu'on sait faire, on n'a pas besoin
%  d'en parler.
% ===================================================================

%%K t10-apar-1 apar
\VUsaut{4.0mm}
\VUtitre{15.0pt}{60}{કોષ્ટક નં. ૧૦}
\VUsaut{3.0mm}
\VUpk{11.4pt}{40}{દરિયો. --- બંદર.}
\VUsaut{3.0mm}

%%K t10-01-1 p
1. --- ઉનાળામાં, જ્યારે કેટલાક પર્વત પર શાંતિ અને ઠંડક શોધે છે, ત્યારે
બીજા દરિયાકાંઠે જવાનું પસંદ કરે છે. ગયા વર્ષે અમે એમ જ કર્યું, કેમ કે
અમને \VUgras{મહાસાગર} \textsuperscript{(1)} ખૂબ ગમે છે.

%%K t10-02-1 p
2. --- મારે માટે તો, પાણીના એ અપાર વિસ્તારને જોવામાં મને ભારે આનંદ થાય
છે, જે \VUgras{ક્ષિતિજ} \textsuperscript{(2)} સુધી ફેલાયેલો છે, અને
લાંબી \VUgras{સફર} માટે ઊપડતાં પ્રચંડ \VUgras{વરાળનાં જહાજ}
\textsuperscript{(3)} જોવામાં, જેમાં અમેરિકા જનારા મુસાફરો ચડે છે.

%%K t10-03-1 p
3. --- એટલે જ મારા પિતા, જે મારો શોખ જાણે છે, રજાઓ ગાળવા માટે હંમેશાં
કોઈ ખૂબ શાંત દરિયાકિનારો પસંદ કરે છે, જે અમારાં મોટાં \VUgras{લશ્કરી
બંદરોમાંના} એકની પાસે આવેલો હોય.

%%K t10-03-2 p
અમારા છેલ્લા મુકામ દરમિયાન \VUgras{કાફલો} પ્રચંડ \VUgras{પાળની}
\textsuperscript{(4)} પાછળ લંગર નાખવા આવ્યો, જે \VUgras{લશ્કરી બંદરને}
\textsuperscript{(5)} બંધ કરે છે અને તેની રક્ષા કરે છે, અને મને મારી
ઇચ્છા પ્રમાણે જુદાં જુદાં \VUgras{લડાયક જહાજ} જોવાની તક મળી : નાની
\VUgras{સબમરીન} \textsuperscript{(10)}, જે ડૂબકી મારીને પાણી નીચે અદૃશ્ય
થઈ જાય છે, અને પ્રચંડ \VUgras{બખ્તરબંધ જહાજ} \textsuperscript{(9)} પણ,
જે આજ સુધી લગભગ માત્ર \VUgras{ટૉર્પિડો જહાજના} \textsuperscript{(8)}
હુમલાથી જ ડરતું હતું.

%%K t10-tit-2 sub
\VUsaut{3.0mm}
\VUpk{10.2pt}{40}{નાનાં જહાજ.}
\VUsaut{3.0mm}

%%K t10-04-1 p
4. --- એ નાનાં જહાજ તો પહેલેથી જ ખૂબ \VUgras{કુશળતાથી} જાડા ધાતુના
બખ્તરનું નબળું બિંદુ શોધી કાઢતાં જાણતાં હતાં, ત્યાં જોખમી
\VUgras{ટૉર્પિડો} ચોંટાડવા માટે, જેનો ધડાકો જહાજનો નાશ કરે છે ; તેમ
છતાં એ સબમરીન કરતાં ઓછાં ડરામણાં હતાં, કેમ કે ટૉર્પિડો-વિરોધી જહાજ
તેમની પાછળ સહેલાઈથી પડી શકે છે.

%%K t10-05-1 p
5. --- મેં ઝડપી બખ્તરવાળા \VUgras{ક્રૂઝર} \textsuperscript{(6)} પણ જોયા,
જે સાચા બખ્તરબંધ જહાજ જેટલા જ લગભગ અભેદ્ય છે, કેટલાંક \VUgras{વાહક
જહાજ} \textsuperscript{(12)}, ત્રણ-ચાર \VUgras{તોપ-હોડી}
\textsuperscript{(7)} અને છેવટે એક \VUgras{ફ્રિગેટ}
\textsuperscript{(11)}. \VUgras{એ કાફલો જ્યારે લંગર ઉપાડવા દાવપેચ કરે છે
ત્યારેનું દૃશ્ય સૌથી રસપ્રદ હોય છે.}

%%K t10-06-1 p
6. --- એ મોટા પોલાદના ઢગલા અને સુંદર \VUgras{ત્રણ ડોલકાઠીવાળાં વહાણ} કે
નમણી \VUgras{સઢવાળી હોડીઓ} \textsuperscript{(13)} વચ્ચે બાંધણીનો કેટલો
ફેર છે — એ હોડીઓ આજે પણ વેપારી કાફલામાં વપરાય છે, અને મેં તેમને પૂરા
સઢે બંદરમાં પેસતાં જોઈ, જ્યારે હું \VUgras{ધક્કા} \textsuperscript{(14)}
પર લટાર મારતો હતો. ખરું જોતાં, \VUgras{પવન} સામો હોય \VUgras{ત્યારે} એ
પોતાના ઘણા ફાયદા ગુમાવે છે. બેસિનમાં પાછા પેસવા માટે તેમને બધા સઢ ઉતારી
નાખવા પડે છે, અને તેમને લેવા આવવા માટે \VUgras{ખેંચનારું જહાજ}
\textsuperscript{(16)} જરૂરી બને છે, જે તેમને સાદી \VUgras{માલવાહક
હોડીઓની} \textsuperscript{(17)} જેમ ઘાટ સુધી ખેંચી લાવે છે.

%%K t10-tit-3 sub
\VUsaut{3.0mm}
\VUpk{10.2pt}{40}{વેપારી બંદર.}
\VUsaut{3.0mm}

%%K t10-07-1 p
7. --- હું જે બંદરની વાત કરું છું તેમાં જે રસપ્રદ છે તે એ કે તેમાં માત્ર
લશ્કરી બંદર જ નથી, જે પ્રચંડ કામ કરીને કૃત્રિમ રીતે ગોઠવાયું છે, પણ એક
અદ્ભુત \VUgras{વેપારી બંદર} પણ છે, જે એક મોટી નદીના \VUgras{નદીમુખમાં}
આવેલું છે.

%%K t10-08-1 p
8. --- જે જમીનની જીભ બે બેસિનને જુદાં પાડે છે તે કિલ્લેબંધ છે અને
દરિયા તરફ આગળ વધતી \VUgras{તોપખાનાં} અને \VUgras{બુરજોની} હારથી
રક્ષાયેલી છે. \VUgras{કોટ} \textsuperscript{(15)} પર લટાર મારવા ખાસ
પરવાનગી જોઈએ. મારા કાકા, જે \VUgras{વહાણવાડાના} \textsuperscript{(18)}
ઇજનેર છે, તેમની મારફતે મને એ સહેલાઈથી મળી ગઈ હતી, અને હું મોટા છાપરાં
નીચે બંધાતાં જહાજ જોવા ગયો.

%%K t10-09-1 p
9. --- મેં તો એક બખ્તરબંધ જહાજનું જળાવતરણ પણ જોયું, અને મને એ રોમાંચક
ક્ષણ યાદ છે જે આખા ટોળાએ અનુભવી, જ્યારે એ પ્રચંડ ઢગલો, પોતાના જ ભરોસે
છોડી દેવાયેલો, પારણા જેવા ચોકઠા પર સરકવા લાગ્યો અને નદીના પાણી પર તરવા
આવ્યો.

%%K t10-10-1 p
10. --- વહાણવાડા સુધી જવા માટે \VUgras{કવાયતના મેદાનમાં}
\textsuperscript{(19)} થઈને જવું પડે છે, જ્યાં છાવણીના સિપાહીઓ રોજ કસરત
કરવા આવે છે ; પછી લોખંડના \VUgras{નાના પુલ} \textsuperscript{(20)} પરથી
પસાર થવાય છે, જે કવાયતના ચોગાનને \VUgras{શસ્ત્રાગાર}
\textsuperscript{(21)} સાથે જોડે છે.

%%K t10-tit-4 sub
\VUsaut{3.0mm}
\VUpk{10.2pt}{40}{શસ્ત્રાગાર અને ગોદીઓ.}
\VUsaut{3.0mm}

%%K t10-11-1 p
11. --- મારા કાકા સાથે મેં એ મોટા મકાનની મુલાકાત લીધી, જે \VUgras{તોપો}
\textsuperscript{(22)}, \VUgras{તોપગોળા} \textsuperscript{(23)} અને
\VUgras{ગોળાથી} ભરેલું છે. મેં \VUgras{ધાતુકામનું કારખાનું}
\textsuperscript{(24)} પણ જોયું, જ્યાં વરાળનાં જહાજનાં \VUgras{બોઇલર}
\textsuperscript{(25)} બને છે.

%%K t10-12-1 p
12. --- એકદમ પાસે જ \VUgras{ગોદીઓ} \textsuperscript{(26)} છે —
સમારકામની ગોદીઓ — અને ખૂબ નોંધપાત્ર \VUgras{તરતી ગોદી}
\textsuperscript{(27)}, જ્યાં હું સમારકામ માટે આવેલા એક જહાજ પર એકદમ
પાસેથી \VUgras{પંખો} \textsuperscript{(28)}, \VUgras{તળિયાનો પાટડો}
\textsuperscript{(29)} અને \VUgras{સુકાન} \textsuperscript{(30)} જોઈ
શક્યો.

%%K t10-13-1 p
13. --- વેપારી બંદર પણ લશ્કરી બંદર જેટલું જ અભ્યાસ કરવા જેવું છે, અને
મેં કેટલાય કલાક ઘાટ પર તાકી રહીને ગાળ્યા, જ્યાં \VUgras{પૈડાંવાળાં
જહાજ} \textsuperscript{(31)} બાંધેલાં હોય છે, જે નદીમાં ઉપર ચડે છે ; અને
\VUgras{ડોલકાઠી ચડાવવાની ક્રેન} \textsuperscript{(32)} કામ કરતી જોઈ, જે
પ્રચંડ માલ પીંછાની જેમ ઊંચકે છે.

%%K t10-tit-5 sub
\VUsaut{4.0mm}
\VUpk{10.2pt}{40}{બંદરનો સુકાની.}
\VUsaut{3.0mm}

%%K t10-14-1 p
14. --- મેં ખાડીમાંથી બહાર નીકળતાં \VUgras{આટલાંટિક પારનાં જહાજ}
\textsuperscript{(34)} જોઈને વાહ પોકારી, તેમના \VUgras{ધ્વજ}
\textsuperscript{(35)} ખુલ્લી હવામાં ફરકતા હતા, અને કુશળ \VUgras{બંદરનો
સુકાની} \textsuperscript{(36)} તેમને દોરતો હતો, જે \VUgras{બોયા}
\textsuperscript{(40)} અને \VUgras{નિશાની થાંભલાથી} ચિહ્નિત
\VUgras{વહાણ-માર્ગે} અદ્ભુત રીતે ચલાવવાનું જાણે છે, \VUgras{સપાટ
તળિયાની હોડીઓથી} \textsuperscript{(41)} બચતાં બચતાં, જે પોતાના
એકમાત્ર ચોરસ સઢથી માંડ માંડ દોરાય છે. મેં \VUgras{આગલા ભાગ}
\textsuperscript{(37)} પર --- વહાણનું નાક --- બીજા વર્ગની
\VUgras{કૅબિન} \textsuperscript{(39)} ઓળખી, અને \VUgras{પાછલા ભાગ}
\textsuperscript{(38)} પર --- વહાણનું પૂછડું --- પહેલા વર્ગના મુસાફરો
માટેના ખંડ.

%%K t10-15-1 p
15. --- ક્યારેક મેં \VUgras{માલવાહક હોડીમાં} \textsuperscript{(42)} માલ
ભરાતો પણ જોયો, જેમાં હમણાં જ \VUgras{વખારનો} \textsuperscript{(43)} અને
\VUgras{દરિયાઈ સ્ટેશનનો} \textsuperscript{(44)} માલ ખડકાયો હતો, જે
\VUgras{ઘાટની રેલલાઇનની} \textsuperscript{(45)} અસંખ્ય ટ્રેનો લાવી હતી.

%%K t10-tit-6 sub
\VUsaut{3.0mm}
\VUpk{10.2pt}{40}{હલેસાંનો આનંદ.}
\VUsaut{3.0mm}

%%K t10-16-1 p
16. --- હું ઘણી વાર \VUgras{ભરતી-ગોદીમાં} \textsuperscript{(53)} હોડી
ચલાવતો, જ્યાં \VUgras{હોડીઓ} \textsuperscript{(46)} આશરો લેવા આવે છે,
અને \VUgras{હલેસું} \textsuperscript{(47)} ચલાવવું મારે માટે આનંદ હતો.
હું એક \VUgras{તૂતક} \textsuperscript{(48)} પર ચડ્યો, જે \VUgras{સઢવાળી
હોડીનું} \textsuperscript{(49)} હતું ; \VUgras{દોરડાં}
\textsuperscript{(52)} વડે \VUgras{ડોલકાઠી} \textsuperscript{(51)} પર
ચડીને \VUgras{આડકાઠી} \textsuperscript{(50)} સુધી પહોંચ્યો ; પછી મેં
\VUgras{હલકી હોડીથી} \textsuperscript{(54)} મારી લટાર ફરી શરૂ કરી, ભારે
\VUgras{માછીમારીની હોડીઓ} \textsuperscript{(55)} વચ્ચે, જ્યાં
\VUgras{જાળ} \textsuperscript{(56)} સુકાય છે.

%%K t10-17-1 p
17. --- \VUgras{ઘાટ} \textsuperscript{(58)} પર મારી સામેથી જાતજાતનો
\VUgras{માલ} \textsuperscript{(59)} પસાર થયો, \VUgras{પેટીઓમાં}
\textsuperscript{(60)} કે \VUgras{ગાંસડીઓમાં} \textsuperscript{(61)}
ભરેલો. એ \VUgras{માલવાહક હોડીઓનો} \VUgras{ભરેલો માલ}
\textsuperscript{(63)} બનતો હતો, અને શક્તિશાળી \VUgras{ક્રેન}
\textsuperscript{(62)} તેને ઉપાડીને \VUgras{ટ્રક} \textsuperscript{(64)}
પર મૂકતી હતી, જ્યારે \VUgras{ટ્રકવાળા} \VUgras{જકાતઘરમાં}
\textsuperscript{(65)} એની જાહેરાત કરીને કર ભરતા હતા.

%%K t10-18-1 p
18. --- છેવટે, જ્યારે હું \VUgras{ફરતા પુલ} \textsuperscript{(66)} આગળ
કે \VUgras{પાણીના દરવાજા} \textsuperscript{(69)} આગળ પહોંચ્યો,
\VUgras{કાંપ ઉલેચવાના યંત્ર} \textsuperscript{(68)} આગળથી પસાર થઈને, જે
\VUgras{નહેરને} \textsuperscript{(67)} સાફ કરે છે, ત્યારે હું
\VUgras{સંકેત-થાંભલા} \textsuperscript{(70)} પાસે ધક્કા પર ઊતર્યો.

%%K t10-19-1 p
19. --- એ ધક્કાની બીજી બાજુએ જ \VUgras{મછવા} \textsuperscript{(71)}
આવીને લાગે છે, જે કાફલાના અફસરોને જમીન પર લાવે છે. ખલાસીઓને તાલમાં
પોતાનાં હલેસાં ઉપાડતા જોવા એ જોવા જેવું દૃશ્ય છે, જ્યારે \VUgras{મોજા}
\textsuperscript{(72)} પર ઊંચકાયેલી હોડી સહેલાઈથી ઊતરવાનાં પગથિયાં પાસે
આવી લાગે છે. આ જગ્યાએથી \VUgras{ભરતીઓટ} અને \VUgras{દરિયાકિનારે}
\textsuperscript{(73)} \VUgras{ભરતી} અને \VUgras{ઓટની} અવરજવર વધારે
સારી રીતે જોઈ શકાય છે.

%%K t10-tit-7 sub
\VUsaut{3.0mm}
\VUpk{10.2pt}{40}{ક્લબઘર.}
\VUsaut{3.0mm}

%%K t10-20-1 p
20. --- ઘરે પાછા ફરવા મેં \VUgras{વીજળીની ટ્રામ} \textsuperscript{(74)}
લીધી, જેનો \VUgras{થોભો} \textsuperscript{(75)} અફસરોની ક્લબ આગળ મુકાયેલા
\VUgras{થાંભલા} પાસે છે.

%%K t10-20-2 p
ક્લબના \VUgras{ઝરૂખા} \textsuperscript{(76)} પરથી, જે \VUgras{લંગર}
\textsuperscript{(77)}, તોપો અને તોપગોળાથી શણગારેલો છે, મેં ઘણી વાર
નૌકા-અફસરોનું ભવ્ય મંડળ જોયું છે : \VUgras{લેફ્ટનન્ટો}
\textsuperscript{(78)}, પોતાની તલવાર સાથે ; \VUgras{તોપખાનાના}
\textsuperscript{(79)} અને \VUgras{પાયદળના} \textsuperscript{(80)}
\VUgras{કપ્તાનો}, પોતાની \VUgras{વાંકી તલવાર} સાથે. તેમની પાછળ એક
\VUgras{ખલાસી} \textsuperscript{(81)} હતો, જે તેમને \VUgras{દૂરબીન}
આપતો, જ્યારે તેઓ દૂર શું બની રહ્યું છે તે ઓળખવા માગતા.

%%K t10-21-1 p
21. --- મેં જુવાન \VUgras{સબ-લેફ્ટનન્ટોને} \textsuperscript{(82)} પણ
જોયા, પોતાના \VUgras{કમાન્ડર} \textsuperscript{(83)} કે
\VUgras{એડમિરલ} \textsuperscript{(84)} પાસેથી સૂચનાઓ લેતા, જ્યારે એક
\VUgras{ક્વાર્ટરમાસ્ટર} \textsuperscript{(85)} એ સૂચનાઓ જહાજ પર લઈ જવા
રાહ જોતો હતો.

%%K t10-22-1 p
22. --- એ ઝરૂખેથી દેખાવ ભવ્ય છે : આખા દરિયાકિનારા પર નજર પડે છે, તેના
\VUgras{તંબૂ} \textsuperscript{(86)} અને \VUgras{કપડાં બદલવાની કૅબિન}
\textsuperscript{(87)} સાથે, અને નહાવાના સમયે \VUgras{નહાનારાઓ}
\textsuperscript{(88)} દેખાય છે, જે \VUgras{જુગારઘર}
\textsuperscript{(89)} આગળ આનંદથી ધમાલ કરે છે.

%%K t10-23-1 p
23. --- કિનારાને છેડે દરિયાકાંઠો એક નાનો \VUgras{દ્વીપકલ્પ}
\textsuperscript{(91)} \VUgras{ખડકાળ} બનાવે છે, જ્યાં \VUgras{મોટાં
મોજાં} \textsuperscript{(92)} આવીને ભાંગે છે અને જેના પર
\VUgras{દીવાદાંડી} \textsuperscript{(93)} બાંધેલી છે, જે રાત્રે
વહાણચાલકોને રસ્તો બતાવે છે. એ દ્વીપકલ્પ જમીન સાથે માત્ર એક ખૂબ સાંકડી
\VUgras{સંયોગીભૂમિથી} \textsuperscript{(90)} જોડાયેલો છે.

%%K t10-tit-8 sub
\VUsaut{3.0mm}
\VUpk{10.2pt}{40}{દરિયાકાંઠાનું સામાન્ય દૃશ્ય.}
\VUsaut{3.0mm}

%%K t10-24-1 p
24. --- \VUgras{ભેખડ} \textsuperscript{(94)} લંબાય છે, દરિયાએ ચીરેલી, જે
તેમાં મરજી મુજબ \VUgras{ખાડીઓ} \textsuperscript{(95)} અને
\VUgras{ભૂશિરોની} હાર દોરે છે અને તેને સૌથી રમણીય બનાવે છે.

%%K t10-24-2 p
\VUgras{ઉચ્ચપ્રદેશ} \textsuperscript{(96)} પર, જે \VUgras{સામુદ્રધુની}
\textsuperscript{(98)} પર પ્રભુત્વ ધરાવે છે, એક ખૂબ મહત્ત્વનો
\VUgras{કિલ્લો} \textsuperscript{(97)} અને થોડા « વિલા » છે.

%%K t10-25-1 p
25. --- એ સામુદ્રધુની ખંડથી એક ખૂબ જોખમી \VUgras{ટાપુને}
\textsuperscript{(99)} જુદો પાડે છે, જેની ચારે બાજુ \VUgras{પાણી નીચેના
ખડક} \textsuperscript{(100)} અને \VUgras{ભાંગતાં મોજાં} છે, જ્યાં
હોડીઓ અને મોટાં જહાજ સુધ્ધાં ક્યારેક ખરાબે ચડે છે. બચાવની તૈયારી અને
\VUgras{બચાવ-હોડીઓના} ઝડપી આગમન છતાં એ \VUgras{વહાણભંગ} ભયાનક હોય છે,
અને \VUgras{સમરાત્રિનાં તોફાનો} વખતે કેટલાંય જહાજ માણસો અને માલ સાથે
ડૂબ્યાં છે.

%%K t10-25-2 p
આટલું પડોશમાં હોવા છતાં \VUgras{એ બંદરે} ખલાસીઓની ભારે અવરજવર રહે છે.

\VUsaut{6.0mm}
\VUfilet{22mm}
